\documentclass[11pt]{article}
\usepackage[margin=25mm]{geometry}
\usepackage{fontspec}
\setmainfont{Latin Modern Roman}
\usepackage{amsmath,amssymb,mathtools}
\usepackage{unicode-math}
\setmathfont{Latin Modern Math}
\usepackage{xurl,hyperref}
\hypersetup{colorlinks=true,urlcolor=blue}
\setcounter{secnumdepth}{0}
\setlength{\parindent}{0pt}
\setlength{\parskip}{5pt}
\setlength{\emergencystretch}{3em}
\newcommand{\sref}[2]{\hyperlink{source:#1:#2}{[S#2]}}
\newcommand{\cataloguescope}[1]{\paragraph{Recorded target.}#1}
\newcommand{\R}{\mathbb R}
\newcommand{\C}{\mathbb C}
\newcommand{\N}{\mathbb N}
\newcommand{\Z}{\mathbb Z}
\newcommand{\Q}{\mathbb Q}
\newcommand{\F}{\mathbb F}
\newcommand{\eps}{\varepsilon}
\begin{document}
\section{First-Principles Derivation of Adaptive-Resolution Molecular Dynamics}\label{statement}

\subsection{Definitions and mathematical statement}

\paragraph{Shared notation.}$\mathbb N=\{1,2,\ldots\}$; the standard number fields have their usual meanings. All additional conventions are stated locally.


\paragraph{Definitions and conventions.} Adaptive-resolution molecular dynamics couples an atomistic region and a coarse-grained region using a switching function $w\in[0,1]$. For molecular centers $X_\alpha,X_\beta$, the report's force interpolation has atomistic weight $w(X_\alpha)w(X_\beta)$ and coarse-grained weight $1-w(X_\alpha)w(X_\beta)$. A potential interpolation generally yields additional derivatives of the spatial weights, hence drift forces. A first-principles derivation would obtain the prescribed interpolation and statistical ensemble from microscopic statistical mechanics rather than choosing them ad hoc.

\paragraph{Source scope and qualifications.} The report supplies an argument against a particular potential-interpolation scheme and leaves a first-principles force derivation open. It does not establish that every conceivable energy-based coarse-graining is impossible; that stronger conclusion is not part of the target.

\paragraph{Catalogue mathematical target.}
\paragraph{Catalogue wording retained for traceability.}
For the adaptive-resolution interpolation formula in molecular dynamics, why is it not naturally formulated using potentials, and can the force-based interpolation formula be formally derived from first principles of statistical physics?

\paragraph{Saved literature context.} The supplied catalogue's assessment dated 2026-08-21 records: The potential-based coupling is ruled out by an overdetermined drift system in the report, while a first-principles derivation of the force interpolation formula is explicitly left as a challenge. The assessment is retained as dated provenance; it is not a new comprehensive literature or open-status certification.
\subsection{Short English statement}
Explain why the adaptive-resolution coupling is force based and derive that force interpolation from statistical mechanics.
\subsection{Research attempt}
\paragraph{Model, date and investigated scope.}
OpenAI GPT-6 Astra (configured model \texttt{gpt-6-astra}), ultra reasoning; generated 11 October 2026. Derives the interpolation drift and an explicit mixed-derivative obstruction for the prescribed force law.

\paragraph{Sources and literature check.}
The complete supplied definition, including source qualifications and the dated literature assessment, was reviewed. The calculation below is an elementary restricted attempt motivated by [S2]; it does not claim novelty or a fresh literature-status determination.

\paragraph{Context and investigated scope.} The complete supplied definition, including source qualifications and the dated literature assessment, was reviewed. The calculation below is an elementary restricted attempt motivated by [S2]; it does not claim novelty or a fresh literature-status determination.

\paragraph{The drift-force calculation.} Let two molecular positions be $x,y$, let $a(x,y)=w(x)w(y)$ and set $U=aU_A(x-y)+(1-a)U_C(x-y)$. Differentiation gives
\[
 -\nabla_xU=aF_A+(1-a)F_C-(\nabla_xa)(U_A-U_C),
\]
where $F_A=-\nabla_xU_A$ and similarly for $F_C$. The desired interpolation omits the last term, which vanishes only under extra restrictions such as constant $a$ or equality of the potentials on the switching region.

\paragraph{An explicit integrability obstruction.} In one dimension choose $U_A(r)=r^2/2$, $U_C=0$ and smooth $w$ taking values in $[0,1]$. The prescribed pair forces are $F_x=-a(x,y)(x-y)$ and $F_y=a(x,y)(x-y)$. A twice differentiable common potential would require $\partial_yF_x=\partial_xF_y$. Subtracting yields
\[
 \partial_yF_x-\partial_xF_y=-(x-y)(\partial_xa+\partial_ya).
\]
Choose a globally smooth $[0,1]$-valued $w$ equal to $z$ on $(1/4,3/4)$. For $x,y$ in that interior switching interval, this is $-(x-y)(x+y)$, nonzero when $x\ne y$. The forces obey action-reaction but are not conservative on that open region.

\paragraph{Source-level limit.} This is a concrete obstruction to the specified position-weighted force law admitting an ordinary two-coordinate potential, not a theorem against every energy-based adaptive model. Deriving its stationary ensemble or compensating reservoir forces from microscopic statistical mechanics is separate. The first-principles derivation requested in the source is not obtained merely by proving nonconservativity; it remains the first substantial gap after the drift calculation.

\paragraph{Verification and adversarial audit.}
Checked drift signs, action-reaction, mixed derivatives and a switching function in [0,1]; avoided excluding all energy-based coarse-graining. Model self-audit only; no independent human proof checking. Prompts A and B were applied to this canonical record and its qualified source scope. Checked drift signs, action-reaction, mixed derivatives and a switching function in [0,1]; avoided excluding all energy-based coarse-graining. Retained unresolved status and zero original-target progress; A statistical-mechanical derivation of the force interpolation and its stationary ensemble is missing.

Coordinator checked the drift-force sign and both mixed derivatives in the switching region. The globally smooth switching function is specified on an interior interval. This additional check was model-assisted, not independent human verification.

\paragraph{Tools, process and contributions.}
Python 3 was used for file preparation and structural checks; no numerical experiment or theorem prover was used. Contribution prompt A was followed by prompt B. GPT-6 Astra with ultra reasoning drafted the research and performed its stated self-audit; any additional coordinator contribution is identified in the verification and audit. A GPT-6-based Codex coordinator assembled the manuscript and checked repository integration; its exact variant and reasoning configuration were not exposed. No human mathematical contribution, independent human proof check or formal proof certification is supplied. Complete derivations and the consulted source references are included here; private reasoning transcripts are not published.

\paragraph{Final status and first remaining gap.}
\textbf{UNRESOLVED}. A statistical-mechanical derivation of the force interpolation and its stationary ensemble is missing.

\section{Completion Estimate}
\noindent\textbf{COMPLETION: 0\%}
The calculation is an auxiliary or restricted attempt and establishes no newly verified portion of the full source target. A statistical-mechanical derivation of the force interpolation and its stationary ensemble is missing.

\subsection{Sources}
\begin{itemize}
\item[\hypertarget{source:statement:1}{S1}] ULAM.AI and the UnsolvedMath Contributors, UnsolvedMath, catalogue entry. Catalogue curation under CC BY 4.0; definitions and English exposition added locally.
\url{https://huggingface.co/datasets/ulamai/UnsolvedMath}.
\item[\hypertarget{source:statement:2}{S2}] Atomistic Models of Materials: Mathematical Challenges (2008), Oberwolfach Reports, DOI 10.4171/owr/2008/21, pp. 1079--1156; official report source and mathematical context.
\url{https://doi.org/10.4171/owr/2008/21}.
\item[\hypertarget{source:statement:3}{S3}] Official EMS Press PDF of the same Oberwolfach report; PDF pages 28--30.
\url{https://ems.press/content/serial-article-files/46165}.
\item[\hypertarget{source:statement:4}{S4}] Delle Site 2007; reference in the saved dated catalogue literature review.
\url{https://doi.org/10.1103/PhysRevE.76.047701}.
\end{itemize}
\label{end:statement}
\end{document}
